\chapterbackmatter{Solutions to Weinberg Exercises}
\ExerciseEditorialNote

\begin{enumerate}
\WeinbergSolution{1}
Introduce the transverse and longitudinal projectors
\[
 P^T_{\mu\nu}(k)=\eta_{\mu\nu}-\frac{k_\mu k_\nu}{k^2},
 \qquad
 P^L_{\mu\nu}(k)=\frac{k_\mu k_\nu}{k^2}.
\]
The inverse free Proca propagator is
\[
 K^{(0)}_{\mu\nu}
 =(k^2+m^2)P^T_{\mu\nu}+m^2P^L_{\mu\nu}.
\]
Lorentz invariance permits the one-particle-irreducible two-point
function to be written
\[
 \Pi^*_{\mu\nu}(k)
 =\Pi_T(k^2)P^T_{\mu\nu}+\Pi_L(k^2)P^L_{\mu\nu}.
\]
Consequently the complete propagator has a transverse denominator
\(k^2+m^2-\Pi_T(k^2)\) and a longitudinal denominator
\(m^2-\Pi_L(k^2)\).  A spin-one particle of physical mass \(m\) and unit
pole residue therefore requires
\[
 \boxed{\Pi_T(-m^2)=0,\qquad \Pi_T'(-m^2)=0.}
\]
The longitudinal denominator must be nonzero at \(k^2=-m^2\); otherwise
an additional scalar state would occur at the same mass.  More generally,
any zero of \(m^2-\Pi_L(k^2)\) represents an extra spin-zero excitation
and must be absent if the field is to describe only the stated particle.

Starting from
\[
 \mathcal L_B=-\frac14v_{B\mu\nu}v_B^{\mu\nu}
 -\frac12m_B^2v_{B\mu}v_B^\mu-V_B(v_B),
 \qquad
 v_{B\mu\nu}=\partial_\mu v_{B\nu}-\partial_\nu v_{B\mu},
\]
put
\[
 v_{B\mu}=\sqrt Z\,v_\mu,\qquad
 m^2=m_B^2+\delta m^2.
\]
Then one convenient split is
\[
 \mathcal L_0=-\frac14v_{\mu\nu}v^{\mu\nu}
 -\frac12m^2v_\mu v^\mu
\]
and
\[
 \mathcal L_1=-\frac14(Z-1)v_{\mu\nu}v^{\mu\nu}
 -\frac12\bigl[(Z-1)m^2-Z\delta m^2\bigr]v_\mu v^\mu
 -V_B(\sqrt Z\,v).
\]
The two local counterterms are chosen so that the transverse loop
self-energy obeys the two boxed on-shell conditions.

\WeinbergSolution{2}
Let \(\Phi\) be a scalar of charge \(q\), and define its complete
propagator and amputated current vertex by
\[
 -i\Delta'(k)\delta^4(k-\ell)
 =\int d^4y\,d^4z\,e^{-ik\cdot y+i\ell\cdot z}
 \bra{\Omega}T\{\Phi(y)\Phi^\dagger(z)\}\ket{\Omega},
\]
\[
 \begin{split}
 &\int d^4x\,d^4y\,d^4z\,
 e^{-ip\cdot x-ik\cdot y+i\ell\cdot z}
 \bra{\Omega}T\{J^\mu(x)\Phi(y)\Phi^\dagger(z)\}\ket{\Omega}\\
 &\qquad=-iq\,\Delta'(k)\Gamma^\mu(k,\ell)
 \Delta'(\ell)\delta^4(p+k-\ell).
 \end{split}
\]
Differentiating the time-ordered product gives contact terms.  Current
conservation and the equal-time charge commutators
\[
 [J^0(\mathbf x,t),\Phi(\mathbf y,t)]
 =-q\Phi(\mathbf y,t)\delta^3(\mathbf x-\mathbf y),
\]
\[
 [J^0(\mathbf x,t),\Phi^\dagger(\mathbf y,t)]
 =q\Phi^\dagger(\mathbf y,t)\delta^3(\mathbf x-\mathbf y)
\]
give
\[
 \partial_\mu^xT\{J^\mu(x)\Phi(y)\Phi^\dagger(z)\}
 =-q\delta^4(x-y)T\{\Phi(y)\Phi^\dagger(z)\}
 +q\delta^4(x-z)T\{\Phi(y)\Phi^\dagger(z)\}.
\]
Fourier transformation and amputation therefore yield the generalized
scalar Ward identity
\[
 \boxed{\;
 (\ell-k)_\mu\Gamma^\mu(k,\ell)
 =i\Delta'^{-1}(k)-i\Delta'^{-1}(\ell).\;}
\]
At zero momentum transfer this becomes the ordinary Ward identity
\[
 \boxed{\;
 \Gamma^\mu(k,k)
 =-i\frac{\partial}{\partial k_\mu}\Delta'^{-1}(k).\;}
\]
For a free scalar \(\Delta'^{-1}(k)=k^2+m^2\), so
\(\Gamma^\mu(k,\ell)=-i(k+\ell)^\mu\), which verifies both formulas and
fixes the convention-dependent factors of \(i\).

\WeinbergSolution{3}
Write \(q=p_2-p_1\), \(P=p_2+p_1\), and suppress the state-normalization
factor.  For states of equal intrinsic parity, parity and Lorentz
covariance allow
\[
 \bra{p_2,\sigma_2}J^\mu(0)\ket{p_1,\sigma_1}
 =\bar u_2\bigl[
 \gamma^\mu F(q^2)+P^\mu G(q^2)+q^\mu H(q^2)
 \bigr]u_1 .
\]
The on-shell equations imply
\[
 \bar u_2\slashed q\,u_1
 =i(m_2-m_1)\bar u_2u_1,\qquad
 q\cdot P=-(m_2^2-m_1^2).
\]
Thus current conservation imposes one relation,
\[
 i(m_2-m_1)F-(m_2^2-m_1^2)G+q^2H=0,
\]
leaving two independent form factors.  A manifestly transverse basis is
\[
 \boxed{\;
 \bar u_2\left[
 \left(\gamma^\mu-\frac{q^\mu\slashed q}{q^2}\right)F_1(q^2)
 +\frac{i\sigma^{\mu\nu}q_\nu}{m_1+m_2}F_2(q^2)
 \right]u_1,\;}
\]
where \(\sigma^{\mu\nu}\equiv\tfrac{i}{2}
[\gamma^\mu,\gamma^\nu]\); harmless factors of \(i\) may instead be
absorbed in the definitions of \(F_1,F_2\).  For unequal masses, regularity
at \(q^2=0\) requires \(F_1(0)=0\).  This is also the statement that the
conserved charge cannot connect distinct mass eigenstates.

If the two states have opposite intrinsic parity, the allowed matrix
element is obtained by inserting \(\gamma_5\):
\[
 \boxed{\;
 \bar u_2\left[
 \left(\gamma^\mu-\frac{q^\mu\slashed q}{q^2}\right)
 \widetilde F_1(q^2)
 +\frac{i\sigma^{\mu\nu}q_\nu}{m_1+m_2}
 \widetilde F_2(q^2)
 \right]\gamma_5u_1.\;}
\]
Equivalently one may start with
\((\gamma^\mu\widetilde F+P^\mu\widetilde G+
q^\mu\widetilde H)\gamma_5\).  In that basis the conservation relation
contains \(m_2+m_1\), because
\(\bar u_2\slashed q\gamma_5u_1
=i(m_2+m_1)\bar u_2\gamma_5u_1\).  Again
\(\widetilde F_1(0)=0\) for a nonsingular transition matrix element.

\WeinbergSolution{4}
Define
\[
 G^{\mu\nu}(q)=\int d^4x\,e^{-iq\cdot x}
 \bra{\Omega}T\{J^\mu(x)J^\nu(0)\}\ket{\Omega}.
\]
Inserting a complete set of momentum eigenstates into the two Wightman
functions shows that the nonlocal part is supported at
\(p^0>0,\ p^2=-\mu^2\).  Lorentz covariance and
\(p_\mu\bra{\Omega}J^\mu(0)\ket{n,p}=0\) force its tensor spectral
density to be transverse:
\[
 \rho^{\mu\nu}(p)
 =\bigl(p^\mu p^\nu-p^2\eta^{\mu\nu}\bigr)\rho(\mu^2).
\]
Time ordering then gives, up to local contact terms,
\[
 \boxed{\;
 G^{\mu\nu}(q)
 =\bigl(q^2\eta^{\mu\nu}-q^\mu q^\nu\bigr)
 \left[
 P(q^2)+\int_0^\infty
 \frac{\rho(\mu^2)\,d\mu^2}
 {q^2+\mu^2-i\epsilon}
 \right].\;}
\]
Here \(P\) is a subtraction polynomial; changing it corresponds to
changing local terms in the product of the two currents.  Such terms are
unavoidable in general because differentiating a time-ordered product can
produce equal-time commutators.

For a genuinely complex current, the spectral function in the displayed
\(JJ\) correlator need not be nonnegative: it is a sum of products
\(\bra{\Omega}J^\mu\ket n\bra nJ^\nu\ket{\Omega}\), not absolute
squares.  The mixed correlator
\(\bra{\Omega}T\{J^\mu(x)J^{\nu\dagger}(0)\}\ket{\Omega}\) has the same
form with a nonnegative transverse spectral measure.  If a stable
spin-one state of mass \(M\) couples to the current, its contribution is
a positive pole term \(Z\delta(\mu^2-M^2)\) in the mixed spectral
function.  Scalar intermediate states are excluded by current
conservation except for possible massless or contact contributions.

\WeinbergSolution{5}
Insert a complete set between \(\psi(x)\) and \(\bar\psi(y)\), use
translation invariance, and organize the resulting spinor matrix by
Lorentz covariance.  If parity is conserved, only a scalar and a vector
Dirac structure occur.  The exact momentum-space propagator therefore has
the form
\[
 \boxed{\;
 S'(p)=\int_0^\infty d\mu^2\,
 \frac{-i\slashed p\,\rho_1(\mu^2)+\rho_2(\mu^2)}
 {p^2+\mu^2-i\epsilon}.\;}
\]
It is often useful to resolve the spectral measure into intermediate
states of the two rest-frame parities:
\[
 \rho_1(\mu^2)=\rho_+(\mu^2)+\rho_-(\mu^2),\qquad
 \rho_2(\mu^2)=\mu\,[\rho_+(\mu^2)-\rho_-(\mu^2)].
\]
Hilbert-space positivity gives
\[
 \rho_\pm(\mu^2)\geq0,
\qquad
 \boxed{\rho_1(\mu^2)\geq0,\quad
 |\rho_2(\mu^2)|\leq\mu\,\rho_1(\mu^2).}
\]
These inequalities follow equally by evaluating the spectral spinor
matrix in the rest frame and requiring its two energy-projector
eigenvalues to be nonnegative.

If \(\psi\) interpolates a stable fermion of mass \(m\), then
\[
 \rho_1(\mu^2)=Z\delta(\mu^2-m^2)+\rho_1^{\,c}(\mu^2),
\qquad
 \rho_2(\mu^2)=Zm\delta(\mu^2-m^2)+\rho_2^{\,c}(\mu^2),
\]
and hence
\[
 S'(p)=Z\frac{-i\slashed p+m}{p^2+m^2-i\epsilon}
 +\text{continuum}.
\]
For a canonically normalized bare field, the equal-time
anticommutator gives the sum rule
\(\int_0^\infty\rho_1(\mu^2)\,d\mu^2=1\), subject to the usual
ultraviolet regulator qualification.

\WeinbergSolution{6}
Assume, for contradiction, that the forward amplitude for a real photon
on a stable target obeys an unsubtracted dispersion relation.  A photon is
its own antiparticle, so crossing makes the spin-averaged forward amplitude
even in the laboratory energy.  The optical theorem and the unsubtracted
relation would then give at zero energy
\[
 \operatorname{Re}f(0)
 =\frac{1}{2\pi^2}\int_0^\infty\sigma_{\gamma\alpha}(E)\,dE.
\]
The integrand is a total cross-section and is nonnegative.  It is not
identically zero for a target that can absorb or scatter photons, so the
right-hand side is positive.

The low-energy theorem gives a contradictory result.  For a target of
mass \(m\) and electric charge \(Qe\), the forward Thomson amplitude is
\[
 f(0)=-\frac{\alpha Q^2}{m}\leq0.
\]
For a neutral target the leading constant is zero, which is still
inconsistent with a strictly positive integral whenever the target has a
nonzero photoabsorption cross-section.  Thus an unsubtracted relation
cannot be valid.

No assumption about the large-energy behavior of either \(f\) or
\(\sigma\) entered the contradiction.  The conclusion is instead that the
contour at infinity cannot be discarded without additional high-energy
information.  At least one subtraction constant---and in the general
forward relation two subtractions---must retain precisely the low-energy
information that the positive absorptive integral cannot reproduce.

\WeinbergSolution{7}
Regard the complete scalar propagator as a function of
\(z=-p^2\).  Microcausality makes it analytic in the complex \(z\)-plane
away from the positive real axis, while unitarity gives the discontinuity
\[
 \operatorname{Disc}\Delta'(z)
 \equiv\Delta'(z+i0)-\Delta'(z-i0)
 =2\pi i\,\rho(z),\qquad \rho(z)\geq0.
\]
Apply Cauchy's theorem to
\[
 \frac{\Delta'(z')}{(z'-z)(z'+\mu_0^2)^n},
\]
where \(n\) is large enough that the circle at infinity vanishes.
Collapsing the contour onto the cut yields
\[
 \Delta'(p)=P_{n-1}(p^2)
 +(-p^2+\mu_0^2)^n
 \int_0^\infty
 \frac{\rho(\mu^2)\,d\mu^2}
 {(\mu^2+\mu_0^2)^n(p^2+\mu^2-i\epsilon)}.
\]
The polynomial \(P_{n-1}\) contains the \(n\) subtraction constants and
represents local contact terms.

For a canonically normalized scalar field, the equal-time commutator fixes
the leading large-\(|p^2|\) behavior and gives
\(\int\rho(\mu^2)d\mu^2=1\).  The unsubtracted form is then
\[
 \boxed{\;
 \Delta'(p)=\int_0^\infty
 \frac{\rho(\mu^2)\,d\mu^2}
 {p^2+\mu^2-i\epsilon}.\;}
\]
A stable particle contributes
\[
 \rho(\mu^2)=Z\delta(\mu^2-m^2)+\sigma(\mu^2),
\qquad \sigma\geq0,
\]
so the isolated pole has residue \(Z>0\), while the continuum begins at
the lightest multiparticle threshold.  This is the
K\"all\'en--Lehmann representation obtained directly from dispersion
theory.

\WeinbergSolution{8}
Let \(f(\omega)\) be the spin-averaged forward amplitude in the electron
rest frame, normalized so that the forward differential cross-section is
\(|f|^2\).  The renormalized low-energy theorem fixes
\[
 f(0)=-\frac{\alpha}{m};
\]
there is no independent order-\(e^4\) correction to this subtraction
constant.  Crossing makes the spin-independent amplitude even in
\(\omega\).  The twice-subtracted dispersion relation and optical theorem
therefore give
\[
 \boxed{\;
 f(\omega)=-\frac{\alpha}{m}
 +\frac{\omega^2}{2\pi^2}\,
 \mathcal P\!\int_0^\infty
 \frac{\sigma_{\rm KN}(E)}{E^2-\omega^2}\,dE
 +\frac{i\omega}{4\pi}\sigma_{\rm KN}(\omega)
 +O(e^6).\;}
\]
Everything after the Thomson term is of order \(e^4\).  The required
absorptive input is the tree-level total Klein--Nishina cross-section.
With \(x=E/m\),
\[
 \begin{split}
 \sigma_{\rm KN}(E)=\frac{2\pi\alpha^2}{m^2}\biggl[
 &\frac{1+x}{x^3}
 \left\{\frac{2x(1+x)}{1+2x}-\ln(1+2x)\right\}\\
 &+\frac{\ln(1+2x)}{2x}
 -\frac{1+3x}{(1+2x)^2}
 \biggr].
 \end{split}
\]
This follows by integrating the differential result of Section 8.7.
The imaginary part in the boxed expression is exactly
\(\operatorname{Im}f=\omega\sigma_{\rm KN}/(4\pi)\), so it satisfies the
optical theorem.  The real part is a single convergent principal-value
integral and therefore completes the order-\(e^4\) calculation without a
direct loop integration.

There is an infrared subtlety at \(\omega=0\).  Because the elastic
electron--photon cut begins there and
\(\sigma_{\rm KN}(E)=\sigma_T[1-2E/m+O(E^2/m^2)]\), the dispersive
integral may not be expanded term by term in \(\omega^2/E^2\).  Its
low-energy boundary value instead contains
\[
 \operatorname{Im}f^{(4)}(\omega)
 =\frac{\sigma_T}{4\pi}\omega+O(\omega^2),
\]
\[
 \operatorname{Re}f^{(4)}(\omega)
 =-\frac{\sigma_T}{\pi^2m}\omega^2
 \ln\frac{m}{|\omega|}+C\omega^2+O(\omega^4\ln|\omega|),
\]
where the finite constant \(C\) is fixed by the full Klein--Nishina
integral.  The logarithm is the expected nonanalyticity from a massless
threshold.
\end{enumerate}
